\chapter{Estimación del circuito cuántico (Ansatz)}
\label{ap:estimacion_ansatz}

En este apéndice se explica, cómo se calcula $\theta$ a partir de los datos y cómo se monta el circuito cuántico (ansatz) usado en QHRP.

\section[Cálculo del vector angular theta]{Cálculo del vector angular $\theta$}
\label{ap:calculo_theta}

Cada dato de entrada $\mathbf{x} = (x_1, \dots, x_P)$ se convierte en un ángulo. Para hacerlo, se compara cada componente con su mínimo y su máximo en la ventana temporal del activo:

\begin{equation}[eq:theta_qhrp]{Definición del ángulo de codificación por característica}
\theta_i =
\begin{cases}
\alpha\,\pi\,\dfrac{x_i - x_i^{\min}}{x_i^{\max} - x_i^{\min}}, & \text{si } x_i^{\max} > x_i^{\min}, \\
0, & \text{si } x_i^{\max} = x_i^{\min},
\end{cases}
\quad i = 1, \dots, P.
\end{equation}

Es decir:
\begin{itemize}
  \item Si la característica $i$ tiene variación, $\theta_i$ indica en qué punto del rango $[x_i^{\min}, x_i^{\max}]$ cae el valor actual.
  \item Si no hay variación (es decir, $x_i^{\max} = x_i^{\min}$), se fija $\theta_i = 0$ para evitar problemas numéricos.
\end{itemize}

El parámetro $\alpha$ controla la amplitud de esos ángulos. Como el factor normalizado está en el intervalo $[0,1]$, se cumple que $\theta_i \in [0, \alpha\pi]$. Por tanto, al usar $\alpha = 2.0$, el rango angular queda en $[0, 2\pi]$.

\subsection[Por qué se fija alpha 2.0]{Por qué se fija $\alpha = 2.0$}

En este trabajo se escoge $\alpha = 2.0$ para que el rango angular de codificación sea exactamente $[0,2\pi]$, es decir, una vuelta completa.

\section{Estructura del ansatz de codificación}
\label{ap:estructura_ansatz}

Una vez calculado $\theta$, el circuito sigue una secuencia de pasos repetidos: se aplican rotaciones para cargar información y, entre medias, se conectan qubits para mezclar esa información. La Tabla~\ref{TB:ANSATZ_BLOQUES} resume el proceso.

\begin{table}[Bloques del ansatz cuántico]{TB:ANSATZ_BLOQUES}{Secuencia de bloques del ansatz usado para codificar una observación de $P$ características en $P$ qubits.}
  \begin{tabular}{p{2.0cm}p{3.5cm}p{8.0cm}}
    \textbf{Bloque} & \textbf{Qué se aplica} & \textbf{Explicación breve} \\
    \hline
    B0 & Inicialización & Se aplica $H$ a todos los qubits para preparar el estado de partida. \\
    B1 & Primera carga & En cada qubit: $T$ y después $R_y(\theta_i)$. \\
    B2 & Primera conexión & Se aplican iSWAP entre pares de qubits definidos. \\
    B3 & Segunda carga & En cada qubit: $T$ y después $R_x(\theta_i)$. \\
    B4 & Segunda conexión & Se repite el bloque iSWAP sobre los mismos pares. \\
    B5 & Tercera carga & En cada qubit: $T$ y después $R_y(\theta_i)$. \\
    B6 & Tercera conexión & Se vuelve a aplicar iSWAP sobre los mismos pares. \\
    B7 & Cierre & En cada qubit: $T$ y $H$ para terminar el circuito. \\
  \end{tabular}
\end{table}

Los pares de qubits que se conectan en la implementación son:

\begin{equation}[eq:pairs_ansatz_qhrp]{Pares de entrelazamiento del ansatz}
\mathcal{E} = \{(0,1), (0,5), (2,4), (0,4), (2,3), (1,2)\}.
\end{equation}

Si se trabaja con menos qubits (por ejemplo, de 6 a 5), un par $(i,j)$ solo se usa cuando ambos índices existen en ese circuito.

\section{Código de referencia en formato de listado}
\label{ap:codigo_theta_ansatz}

Para dejar claro de dónde sale todo, los Códigos~\ref{COD:THETA_QHRP} y~\ref{COD:ANSATZ_QHRP} muestran líneas del archivo original del proyecto (\texttt{src/qhrp\_figure3\_4/src/quantum\_hrp.py}): primero el cálculo de $\theta$ y después la construcción del circuito.

\PythonCode[COD:THETA_QHRP]{Calculo de theta en QHRP.}{Fragmento del archivo original con el cálculo de los ángulos theta.}{../../src/qhrp_figure3_4/src/quantum_hrp.py}{62}{70}{62}

\PythonCode[COD:ANSATZ_QHRP]{Construccion del ansatz en QHRP.}{Fragmento del archivo original con las rotaciones, las conexiones iSWAP y el estado final.}{../../src/qhrp_figure3_4/src/quantum_hrp.py}{71}{107}{71}
